\section{LLM-Assisted Assembly Transformation}\label{sec:llm}
\subsection{What the Model Is Asked to Change}
In this mode the LLM receives assembly and proposes a new instruction
sequence intended to preserve behavior. It supports the EQ, NOP, BCF,
and IBP families introduced in Section~\ref{sec:background}; data
classification instead returns only locations. PII and wrapping remain
native operations because their layout and encoding constraints are not
implemented in the LLM interface's requests and validators. This is an
implementation boundary, not an intrinsic limitation of language models.

Zone prepares an assembly listing with identified data retained separately;
identification itself may use an LLM. Symbolic labels identify locations
independently of changing byte offsets. The model returns a complete listing,
which Zone assembles, restores with the retained data, and validates.
Keeping exact data bytes outside generation avoids requiring the model to
reproduce unchanged content and separates instruction rewriting from data
preservation. An end-to-end alternative is not evaluated. The original binary, prepared listing, and final
binary are different representations: internal LLM similarity compares
prepared and returned listings, excluding changes made during preparation.

The requested model is \texttt{deepseek-flash}. Provider documentation
accessed on 15 November 2026 maps this alias to DeepSeek-V4.1-Flash~\cite{deepseek_models}.
The experiment records retain the alias, not an independently verified
provider snapshot for each call.
Both LLM roles request temperature 0.2, with no sampling seed or explicit
top-p. Repeated requests are not treated as deterministic reproductions.

\subsection{Prompting Conditions and Repeated Attempts}
A system message defines the task, preservation priority, and response
format. The user prompt specifies the architecture, technique-specific
rules, common constraints on machine state, control flow and data,
input-specific notes, and the prepared listing. Uncertain transformations
must remain unapplied; the response must contain only a complete assembly
listing. Both conditions share this structured specification:
\emph{zero-shot} supplies no example, whereas \emph{two-shot} adds two
illustrative input--output examples for the selected technique. The comparison
therefore isolates the addition of demonstrations to the prompt. Shots count
examples, not attempts. Each input, technique, and shot setting defines a
\emph{cell}.

Validation feedback returns a rejected answer and its diagnostic---for
example, a missing label or assembly failure---and requests a correction.
It contains neither a functional success signal nor a scanner verdict.
The controlled benchmark disables correction and makes three independent
single-generation requests per cell, retaining failures. The original-corpus
campaign targets one output per cell, with some retries and failed
restarts in its logs. Its results are cell outcomes rather than uniform
first-generation success rates (Section~\ref{sec:methodology}).

\subsection{From a Proposed Listing to an Assessed Output}
Validation checks whether the response can be assembled and decoded,
whether tracked control-flow targets remain consistent, whether preserved
data survives, and whether the requested transformation's structural
requirements are met. A response can fail at any of these stages, or be
truncated before it contains a complete listing. Passing them means that
Zone accepts a constructed output, which then proceeds to the execution test.

An external \emph{test harness} launches the output and records its result.
It compares the complete expected buffer for controlled programs and
checks the class-specific observable for original payloads
(Section~\ref{sec:methodology}).


Saved records distinguish requests, accepted outputs, subsequent execution
results, time, and estimated cost.
